% =========================================================
%  TEMPLATE 3 — NAVY & GOLD FORMAL ACADEMIC TEMPLATE
%  Dark navy sidebar title page, gold accents, formal serif
% =========================================================
\documentclass[11pt,a4paper]{article}

\usepackage[margin=1in]{geometry}
\usepackage[utf8]{inputenc}
\usepackage[T1]{fontenc}
\usepackage[expansion=false]{microtype}

\usepackage[dvipsnames,table]{xcolor}
\definecolor{navy}{HTML}{0B1F3A}
\definecolor{gold}{HTML}{C9A227}
\definecolor{darktext}{HTML}{1B1B2F}
\definecolor{codebg}{HTML}{F8F7F2}
\definecolor{coderule}{HTML}{E0DCC8}

\usepackage{titlesec}
\titleformat{\section}
  {\normalfont\Large\bfseries\color{navy}}
  {\thesection}{0.6em}{}[{\color{gold}\titlerule[1.4pt]}]
\titleformat{\subsection}
  {\normalfont\large\bfseries\color{navy}}
  {\thesubsection}{0.6em}{}
\titleformat{\subsubsection}
  {\normalfont\normalsize\bfseries\color{gold!70!black}}
  {\thesubsubsection}{0.6em}{}
\titlespacing*{\section}{0pt}{1.4em}{0.8em}

\usepackage{fancyhdr}
\pagestyle{fancy}
\fancyhf{}
\renewcommand{\headrulewidth}{0.8pt}
\renewcommand{\footrulewidth}{0pt}
\fancyhead[L]{\small\color{navy}\bfseries \assignmentTitle}
\fancyhead[R]{\small\color{gold!70!black}\bfseries \courseCode}
\fancyfoot[C]{\small\color{navy}\thepage}
\fancypagestyle{plain}{\fancyhf{}\renewcommand{\headrulewidth}{0pt}\fancyfoot[C]{\small\color{navy}\thepage}}

\usepackage{graphicx}
\usepackage{tikz}
\usepackage{amsmath,amssymb}
\usepackage{enumitem}
\usepackage{booktabs}
\usepackage{array}
\usepackage[most]{tcolorbox}
\tcbuselibrary{skins,breakable}
\usepackage{hyperref}
\hypersetup{colorlinks=true,linkcolor=navy,citecolor=gold!70!black,urlcolor=gold!70!black,pdfborder={0 0 0}}
\usepackage{caption}
\captionsetup{font=small,labelfont={bf,color=navy}}

\usepackage{listings}
\lstdefinestyle{navystyle}{
  backgroundcolor=\color{codebg},
  rulecolor=\color{coderule},
  frame=single,
  framesep=6pt,
  basicstyle=\ttfamily\footnotesize,
  keywordstyle=\color{navy}\bfseries,
  commentstyle=\color{gray}\itshape,
  stringstyle=\color{gold!60!black},
  numberstyle=\tiny\color{gray},
  numbers=left,
  numbersep=8pt,
  breaklines=true,
  showstringspaces=false,
  captionpos=b
}
\lstset{style=navystyle}

\newtcolorbox{experimentbox}[1]{
  breakable, colback=codebg, colframe=navy,
  colbacktitle=navy, coltitle=gold,
  fonttitle=\bfseries, title=#1,
  boxrule=0.9pt, arc=1mm,
  left=6pt,right=6pt,top=6pt,bottom=6pt
}
\newcounter{expcounter}
\newenvironment{experiment}[2]{%
  \refstepcounter{expcounter}%
  \begin{experimentbox}{Experiment \theexpcounter: #1}
  \textbf{\color{navy}Aim:} #2 \par\vspace{4pt}
}{\end{experimentbox}}

\newcommand{\assignmentTitle}{Assignment Title Goes Here}
\newcommand{\courseCode}{Course Code}
\newcommand{\courseName}{Full Course Name}
\newcommand{\studentName}{Student Name}
\newcommand{\studentID}{Roll / Registration Number}
\newcommand{\deptName}{Department Name}
\newcommand{\instituteName}{Institute / University Name}
\newcommand{\instructorName}{Instructor Name}
\newcommand{\submissionDate}{\today}

\begin{document}

% =========================================================
%  TITLE PAGE — sidebar layout
% =========================================================
\begin{titlepage}
\newgeometry{left=0cm,right=0cm,top=0cm,bottom=0cm}
\thispagestyle{empty}
\begin{tikzpicture}[remember picture,overlay]
  % sidebar background
  \fill[navy] (current page.north west) rectangle ([xshift=4.5cm]current page.south west);
  \draw[gold, line width=2pt] ([xshift=4.5cm]current page.north west) -- ([xshift=4.5cm]current page.south west);

  % rotated institute name in sidebar
  \node[gold, rotate=90, font=\Large\bfseries] at ([xshift=2.25cm]current page.west |- current page.center) {\MakeUppercase{\instituteName}};

  % main content column, anchored top-left, fixed x, flowing text block
  \node[anchor=north west, text width=10.5cm, align=left, inner sep=0pt]
    at ([xshift=5.7cm,yshift=-3.2cm]current page.north west) {%
    {\color{navy}\small\bfseries \MakeUppercase{\deptName}}\\[0.3cm]
    {\color{gold!70!black}\large\itshape \courseCode \; --- \; \courseName}\\[1.4cm]
    {\color{navy}\Huge\bfseries \assignmentTitle}\\[0.3cm]
    {\color{gold}\rule{5cm}{2pt}}\\[1.4cm]
    {\large\color{navy}\itshape Submitted in partial fulfillment of the requirements of \courseCode}\\[1.6cm]
    \renewcommand{\arraystretch}{1.7}%
    \begin{tabular}{@{}l l@{}}
    {\color{gold!70!black}\bfseries Submitted by:} & \studentName \\
    {\color{gold!70!black}\bfseries Student ID:}   & \studentID \\
    {\color{gold!70!black}\bfseries Submitted to:} & \instructorName \\
    {\color{gold!70!black}\bfseries Date:}         & \submissionDate \\
    \end{tabular}
  };
\end{tikzpicture}
\end{titlepage}
\restoregeometry

% =========================================================
\tableofcontents
\newpage

\section{Introduction}
Write a brief introduction to the assignment or lab manual here: objective, background,
and scope of the work.

\subsection{Objective}
\begin{itemize}[leftmargin=1.4em]
  \item Objective one.
  \item Objective two.
  \item Objective three.
\end{itemize}

\section{Theory / Background}
Provide the theoretical background, formulas, or concepts relevant to the assignment.

\begin{equation}
  E = mc^2
\end{equation}

\section{Lab Experiments}

Duplicate the \texttt{experiment} block below for every experiment; numbering
increments automatically.

\begin{experiment}{Linear Search Implementation}{To implement and analyze the linear search algorithm.}

\textbf{Procedure:}
\begin{enumerate}[leftmargin=1.4em]
  \item Take an array and a target element as input.
  \item Traverse the array sequentially comparing each element to the target.
  \item Return the index if found, else return $-1$.
\end{enumerate}

\textbf{\color{navy}Source Code:}
\begin{lstlisting}[language=Python, caption={Linear Search in Python}]
def linear_search(arr, target):
    for i, value in enumerate(arr):
        if value == target:
            return i
    return -1

if __name__ == "__main__":
    data = [4, 2, 9, 7, 5]
    result = linear_search(data, 7)
    print(f"Element found at index: {result}")
\end{lstlisting}

\textbf{\color{navy}Output / Observation:}\\
\texttt{Element found at index: 3}

\textbf{\color{navy}Conclusion:}\\
The linear search algorithm successfully locates the target element in
$O(n)$ time complexity in the worst case.

\end{experiment}

\vspace{0.6cm}

\begin{experiment}{Sum of Two Numbers (C Program)}{To write a C program that reads two integers and prints their sum.}

\textbf{\color{navy}Source Code:}
\begin{lstlisting}[language=C, caption={Sum of Two Numbers in C}]
#include <stdio.h>

int main(void) {
    int a, b, sum;
    printf("Enter two integers: ");
    scanf("%d %d", &a, &b);
    sum = a + b;
    printf("Sum = %d\n", sum);
    return 0;
}
\end{lstlisting}

\textbf{\color{navy}Output / Observation:}\\
\texttt{Enter two integers: 5 7}\\
\texttt{Sum = 12}

\end{experiment}

\section{Results and Analysis}
\begin{table}[h]
\centering
\caption{Sample Result Table}
\begin{tabular}{@{}lcc@{}}
\toprule
\textbf{Test Case} & \textbf{Input} & \textbf{Output} \\
\midrule
1 & [4,2,9,7,5], target=7 & Index 3 \\
2 & [1,2,3], target=5     & Not Found \\
\bottomrule
\end{tabular}
\end{table}

\section{Conclusion}
Summarize what was learned or achieved through this assignment / lab exercise.

\section*{References}
\addcontentsline{toc}{section}{References}
\begin{enumerate}[leftmargin=1.4em]
  \item Author, \textit{Title of Book/Resource}, Publisher, Year.
  \item Online resource, \url{https://example.com}, accessed \today.
\end{enumerate}

\end{document}
